\chapter{Products and nonlinear jumps in one-dimensional BV}\label{sec:bv-interval-jump-calculus}
Let $I=(a,b)$ be bounded. We use the right-continuous representatives and left limits from Theorem \ref{thm:interval-bv-representative}. The measure framework is that of \cite{Sim18}, Chapter 2, Section 2; the partition arguments below give explicit one-dimensional proofs of the calculus formulas. At a jump, evaluating an ordinary derivative at only one side generally gives the wrong coefficient.

\begin{theorem}[The divided-difference chain rule on an interval]\label{thm:interval-bv-chain-rule}
Let $u\in BV(I)$ and $\Phi\in C^1(\mathbb R)$. Then $\Phi(u)\in BV(I)$ and
\[
 D(\Phi(u))=K_\Phi(u^-,u^+)\,Du,
 \qquad K_\Phi(s,t)=\int_0^1\Phi'((1-\theta)s+\theta t)\,d\theta.
\]
Equivalently $K_\Phi(s,t)=(\Phi(t)-\Phi(s))/(t-s)$ for $s\ne t$, and $K_\Phi(s,s)=\Phi'(s)$. In particular, at a jump point $x$ the new atom has mass $\Phi(u^+(x))-\Phi(u^-(x))$. Away from the countable jump set the coefficient equals $\Phi'(u^+(x))$.
We prove this one-dimensional formula directly from the measure formulation of BV in \cite{Sim18}, including its jump and singular-continuous parts.
\end{theorem}
\begin{proof}
The regulated representative of $u$ is bounded. On an interval containing its range, $\Phi'$ is bounded; extend the restriction of $\Phi$ to a globally Lipschitz function. Proposition \ref{prop:bv-lipschitz-composition}, together with boundedness of $I$ and the composition, gives $\Phi(u)\in BV(I)$. Its right-continuous representative is $\Phi(u^+)$. The integral defining $K_\Phi$ shows that it is continuous and bounded on the compact square containing all one-sided values.

Fix $a<s<t<b$ and take partitions $s=x_0<\cdots<x_N=t$ with mesh tending to zero. Write $u_i=u^+(x_i)$ and $\mu=Du$. Telescoping and the interval representation give
\[
 \Phi(u^+(t))-\Phi(u^+(s))
 =\sum_iK_\Phi(u_{i-1},u_i)\mu((x_{i-1},x_i]).
\]
Regard the coefficients as a bounded step function on $(s,t]$. For each fixed $r$ in that interval, the upper endpoint of its half-open cell tends to $r$ from the right or equals $r$, and the lower endpoint tends to $r$ strictly from the left. Therefore the corresponding values converge to $u^+(r)$ and $u^-(r)$. Continuity of $K_\Phi$ and dominated convergence with respect to the finite measure $|\mu|$ yield
\[
 \Phi(u^+(t))-\Phi(u^+(s))
 =\int_{(s,t]}K_\Phi(u^-,u^+)\,d\mu.
\]
The left side is $D(\Phi(u))((s,t])$ by Theorem \ref{thm:interval-bv-representative}. Equality on these half-open intervals determines the finite signed measure on $I$, by exhaustion and the uniqueness of a measure on its interval generators. Multiplying the coefficient at an atom by $Du(\{x\})=u^+(x)-u^-(x)$ gives the asserted jump mass.
\end{proof}

\begin{proposition}[The midpoint product rule]\label{prop:interval-bv-product-rule}
For $u,v\in BV(I)$ define $\overline u=(u^-+u^+)/2$ and $\overline v=(v^-+v^+)/2$. Then $uv\in BV(I)$ and
\[
 D(uv)=\overline u\,Dv+\overline v\,Du.
\]
One also has the asymmetric formula $D(uv)=u^+Dv+v^-Du$; the symmetric formula makes the role of simultaneous jumps explicit.
\end{proposition}
\begin{proof}
Boundedness of the representatives implies
\[
 \operatorname{Var}(u^+v^+;I)
 \leq\|u^+\|_\infty\operatorname{Var}(v^+;I)
 +\|v^+\|_\infty\operatorname{Var}(u^+;I).
\]
Alternatively, membership in $BV$ follows from $uv=((u+v)^2-(u-v)^2)/4$ and Theorem \ref{thm:interval-bv-chain-rule}. Apply that theorem to the squares: $D(w^2)=(w^-+w^+)Dw=2\overline w\,Dw$. Linearity of $D$ and of the midpoint representative gives the symmetric formula.

To compare the asymmetric formula, their difference is
\[
 \tfrac12(u^+-u^-)Dv-\tfrac12(v^+-v^-)Du.
\]
Both coefficient functions vanish off their countable jump sets. A nonatomic part of either finite measure gives zero mass to that countable union. At a point $x$, the atom of the first term is $\tfrac12\Delta u(x)\Delta v(x)$, exactly the atom of the second. Thus the measures cancel. This verification includes common jump points and does not assume either function continuous.
\end{proof}

\begin{corollary}[Stieltjes integration by parts with traces]\label{cor:interval-bv-stieltjes-parts}
For $u,v\in BV(a,b)$,
\[
 \int_{(a,b)}\overline u\,dDv+\int_{(a,b)}\overline v\,dDu
 =u(b-)v(b-)-u(a+)v(a+).
\]
In particular, for $\phi\in C^1([a,b])$,
\[
 \int_{(a,b)}\phi\,dDu+\int_a^bu\phi'\,dx
 =u(b-)\phi(b)-u(a+)\phi(a).
\]
\end{corollary}
\begin{proof}
The total mass of $Dw$ for an interval BV function is $w(b-)-w(a+)$ by its cumulative representation. Apply this to $w=uv$ in Proposition \ref{prop:interval-bv-product-rule}. The bounded midpoint representatives are Borel, so both integrals against finite signed measures are well defined. For smooth $\phi$, its two one-sided representatives agree, $D\phi=\phi'\mathcal L^1$, and $\overline u=u$ almost everywhere for Lebesgue measure. This gives the second formula with the stated signs.
\end{proof}

\begin{example}[Squaring a singular-continuous function]\label{ex:cantor-square-derivative}
Let $C$ be the Cantor function on $(0,1)$. Then
\[
 D(C^2)=2C\,\mu_C,
\]
a singular-continuous positive measure of total mass one. In particular $\int C\,d\mu_C=1/2$.
\end{example}
\begin{proof}
The function $C$ is continuous, with $DC=\mu_C$ by Example \ref{ex:cantor-function-bv}. Theorem \ref{thm:interval-bv-chain-rule} therefore reduces to the coefficient $2C$ everywhere. The derivative is singular because it is supported on the Cantor set, and nonatomic because $\mu_C$ is nonatomic. It is positive since $C\geq0$. Its mass equals $C(1-)^2-C(0+)^2=1$ by the trace identity, yielding the integral. This also provides the variation of $C^2$ without replacing its derivative by the almost-everywhere classical derivative, which vanishes off the Cantor set.
\end{proof}

\begin{example}[Why a one-sided derivative is wrong at a jump]\label{ex:nonlinear-bv-jump-coefficient}
For $0<c<1$, set $u=\chi_{[c,1)}$ on $(0,1)$. Then $Du=\delta_c$ and $D(u^k)=\delta_c$ for every integer $k\geq1$. For $k=2$, the formula $2u^+Du$ would instead give $2\delta_c$, and $2u^-Du$ would give zero. For $k=3$, evaluating the ordinary derivative at the midpoint would give $\tfrac34\delta_c$, also incorrect.
\end{example}
\begin{proof}
The function takes only the values zero and one, so $u^k=u$. Its single jump is from zero to one. The chain coefficient for $\Phi(s)=s^k$ is the divided difference $K_\Phi(0,1)=1$, as in Theorem \ref{thm:interval-bv-chain-rule}. For a square this happens to equal $2(0+1)/2$, explaining the midpoint product rule. For a cube the value of $\Phi'$ at $1/2$ is $3/4$, whereas its average along the whole jump is $\int_0^13\theta^2\,d\theta=1$. Thus the general nonlinear rule needs that average, not a single chosen value of the classical derivative.
\end{proof}
